\documentclass[10pt,a4paper]{article}
\usepackage[T1]{fontenc}
\usepackage[utf8]{inputenc}
\usepackage{mathptmx}
\usepackage[a4paper,top=1.2cm,bottom=1.2cm,left=1.45cm,right=1.45cm]{geometry}
\usepackage{enumitem}
\usepackage{titlesec}
\usepackage[hidelinks]{hyperref}
\usepackage{microtype}

% Machine-readable text for applicant tracking systems
\input{glyphtounicode}
\pdfgentounicode=1

\hypersetup{
  pdftitle={Cedric Antony Joseph - Resume},
  pdfauthor={Cedric Antony Joseph},
  pdfsubject={Resume},
  pdfkeywords={Computer Science, Cybersecurity, Linux, Artificial Intelligence, Machine Learning, Python}
}

\pagestyle{empty}
\setlength{\parindent}{0pt}
\setlength{\tabcolsep}{0pt}
\raggedbottom
\raggedright

\titleformat{\section}{\normalsize\bfseries\uppercase}{}{0pt}{}[\vspace{-5pt}\rule{\linewidth}{0.5pt}]
\titlespacing*{\section}{0pt}{7pt}{4pt}

\setlist[itemize]{leftmargin=1.3em,itemsep=1.5pt,topsep=2pt,parsep=0pt,partopsep=0pt,label=\textbullet}

% Organisation line: bold left, date right
\newcommand{\entry}[2]{\textbf{#1}\hfill\textbf{#2}\par}
% Role line: italic left, italic right
\newcommand{\role}[2]{\textit{#1}\hfill\textit{#2}\par}
% Project line: bold name, stack, status right
\newcommand{\project}[3]{\textbf{#1}~\textbar~\textit{#2}\hfill\textbf{#3}\par}

% The public copy carries no phone number. For copies you send directly, build with:
%   pdflatex "\def\phoneline{+00 00000 00000 \textbar{} }\documentclass[10pt,a4paper]{article}
\usepackage[T1]{fontenc}
\usepackage[utf8]{inputenc}
\usepackage{mathptmx}
\usepackage[a4paper,top=1.2cm,bottom=1.2cm,left=1.45cm,right=1.45cm]{geometry}
\usepackage{enumitem}
\usepackage{titlesec}
\usepackage[hidelinks]{hyperref}
\usepackage{microtype}

% Machine-readable text for applicant tracking systems
\input{glyphtounicode}
\pdfgentounicode=1

\hypersetup{
  pdftitle={Cedric Antony Joseph - Resume},
  pdfauthor={Cedric Antony Joseph},
  pdfsubject={Resume},
  pdfkeywords={Computer Science, Cybersecurity, Linux, Artificial Intelligence, Machine Learning, Python}
}

\pagestyle{empty}
\setlength{\parindent}{0pt}
\setlength{\tabcolsep}{0pt}
\raggedbottom
\raggedright

\titleformat{\section}{\normalsize\bfseries\uppercase}{}{0pt}{}[\vspace{-5pt}\rule{\linewidth}{0.5pt}]
\titlespacing*{\section}{0pt}{7pt}{4pt}

\setlist[itemize]{leftmargin=1.3em,itemsep=1.5pt,topsep=2pt,parsep=0pt,partopsep=0pt,label=\textbullet}

% Organisation line: bold left, date right
\newcommand{\entry}[2]{\textbf{#1}\hfill\textbf{#2}\par}
% Role line: italic left, italic right
\newcommand{\role}[2]{\textit{#1}\hfill\textit{#2}\par}
% Project line: bold name, stack, status right
\newcommand{\project}[3]{\textbf{#1}~\textbar~\textit{#2}\hfill\textbf{#3}\par}

% The public copy carries no phone number. For copies you send directly, build with:
%   pdflatex "\def\phoneline{+00 00000 00000 \textbar{} }\documentclass[10pt,a4paper]{article}
\usepackage[T1]{fontenc}
\usepackage[utf8]{inputenc}
\usepackage{mathptmx}
\usepackage[a4paper,top=1.2cm,bottom=1.2cm,left=1.45cm,right=1.45cm]{geometry}
\usepackage{enumitem}
\usepackage{titlesec}
\usepackage[hidelinks]{hyperref}
\usepackage{microtype}

% Machine-readable text for applicant tracking systems
\input{glyphtounicode}
\pdfgentounicode=1

\hypersetup{
  pdftitle={Cedric Antony Joseph - Resume},
  pdfauthor={Cedric Antony Joseph},
  pdfsubject={Resume},
  pdfkeywords={Computer Science, Cybersecurity, Linux, Artificial Intelligence, Machine Learning, Python}
}

\pagestyle{empty}
\setlength{\parindent}{0pt}
\setlength{\tabcolsep}{0pt}
\raggedbottom
\raggedright

\titleformat{\section}{\normalsize\bfseries\uppercase}{}{0pt}{}[\vspace{-5pt}\rule{\linewidth}{0.5pt}]
\titlespacing*{\section}{0pt}{7pt}{4pt}

\setlist[itemize]{leftmargin=1.3em,itemsep=1.5pt,topsep=2pt,parsep=0pt,partopsep=0pt,label=\textbullet}

% Organisation line: bold left, date right
\newcommand{\entry}[2]{\textbf{#1}\hfill\textbf{#2}\par}
% Role line: italic left, italic right
\newcommand{\role}[2]{\textit{#1}\hfill\textit{#2}\par}
% Project line: bold name, stack, status right
\newcommand{\project}[3]{\textbf{#1}~\textbar~\textit{#2}\hfill\textbf{#3}\par}

% The public copy carries no phone number. For copies you send directly, build with:
%   pdflatex "\def\phoneline{+00 00000 00000 \textbar{} }\documentclass[10pt,a4paper]{article}
\usepackage[T1]{fontenc}
\usepackage[utf8]{inputenc}
\usepackage{mathptmx}
\usepackage[a4paper,top=1.2cm,bottom=1.2cm,left=1.45cm,right=1.45cm]{geometry}
\usepackage{enumitem}
\usepackage{titlesec}
\usepackage[hidelinks]{hyperref}
\usepackage{microtype}

% Machine-readable text for applicant tracking systems
\input{glyphtounicode}
\pdfgentounicode=1

\hypersetup{
  pdftitle={Cedric Antony Joseph - Resume},
  pdfauthor={Cedric Antony Joseph},
  pdfsubject={Resume},
  pdfkeywords={Computer Science, Cybersecurity, Linux, Artificial Intelligence, Machine Learning, Python}
}

\pagestyle{empty}
\setlength{\parindent}{0pt}
\setlength{\tabcolsep}{0pt}
\raggedbottom
\raggedright

\titleformat{\section}{\normalsize\bfseries\uppercase}{}{0pt}{}[\vspace{-5pt}\rule{\linewidth}{0.5pt}]
\titlespacing*{\section}{0pt}{7pt}{4pt}

\setlist[itemize]{leftmargin=1.3em,itemsep=1.5pt,topsep=2pt,parsep=0pt,partopsep=0pt,label=\textbullet}

% Organisation line: bold left, date right
\newcommand{\entry}[2]{\textbf{#1}\hfill\textbf{#2}\par}
% Role line: italic left, italic right
\newcommand{\role}[2]{\textit{#1}\hfill\textit{#2}\par}
% Project line: bold name, stack, status right
\newcommand{\project}[3]{\textbf{#1}~\textbar~\textit{#2}\hfill\textbf{#3}\par}

% The public copy carries no phone number. For copies you send directly, build with:
%   pdflatex "\def\phoneline{+00 00000 00000 \textbar{} }\input{resume.tex}"
\providecommand{\phoneline}{}

\begin{document}

\begin{center}
  {\LARGE\bfseries Cedric Antony Joseph}\\[5pt]
  Chennai, India \textbar{} \phoneline
  \href{mailto:cedricaj31@gmail.com}{cedricaj31@gmail.com}\\[2pt]
  \href{https://www.linkedin.com/in/cedric-antony-joseph-5b68b4224/}{linkedin.com/in/cedric-antony-joseph-5b68b4224} \textbar{}
  \href{https://github.com/Divide317}{github.com/Divide317}
\end{center}
\vspace{-6pt}

\section{Summary}
Computer science undergraduate moving from AI and machine learning into cybersecurity and systems. Building detection engines at three layers of the stack: kernel (Linux rootkit detector), device (signed sensor telemetry) and ledger (money-laundering engine).

\section{Education}
\entry{Vellore Institute of Technology (VIT), Chennai}{Expected June 2028}
\role{B.Tech in Computer Science and Engineering (Specialisation in Artificial Intelligence and Machine Learning)}{Chennai, India}
\begin{itemize}
  \item \textbf{Relevant Coursework:} Operating Systems, Computer Networks, Data Structures and Algorithms, Database Management Systems, Machine Learning, Probability and Statistics
  \item \textbf{Certification in progress:} Google Cybersecurity Professional Certificate (Coursera)
\end{itemize}

\section{Technical Skills}
\textbf{Languages:} C, C++, Python, Java, SQL, JavaScript\par
\vspace{1.5pt}
\textbf{Systems and Security:} Linux internals (/proc, syscalls), rootkit detection, MITRE ATT\&CK, threat modelling, Kali Linux\par
\vspace{1.5pt}
\textbf{Security Operations (certificate in progress):} SIEM tools, intrusion detection systems, packet analysis, incident response\par
\vspace{1.5pt}
\textbf{Testing and Tooling:} Git, GitHub Actions, AddressSanitizer, UBSan, Valgrind, clang-tidy, pytest, Docker\par
\vspace{1.5pt}
\textbf{Machine Learning and Backend:} PyTorch, TensorFlow, scikit-learn, OpenCV, ONNX, FastAPI, PostgreSQL, MongoDB, AWS\par

\section{Experience}
\entry{Superposition Technologies}{One-month internship}
\role{AI Intern}{}
\begin{itemize}
  \item Contributed to an AI-driven platform that routes inbound calls to the right recipient and schedules recurring weekly meetings without manual coordination
\end{itemize}

\section{Projects}
\project{Crossview}{C, Python, GitHub Actions}{In progress}
\begin{itemize}
  \item Building a Linux rootkit detector: a C library asks the kernel about processes, modules, ports and preload hooks through independent interfaces (libc, raw syscalls, direct probes) and flags where the answers disagree
  \item Localises each mismatch to user-space or kernel-level hiding and maps it to a MITRE ATT\&CK technique; Python command-line tool over a ctypes interface with JSON reports
  \item CI runs the tests under AddressSanitizer, UBSan and Valgrind with clang-tidy and pytest; published threat model and limitations
\end{itemize}
\vspace{2pt}
\project{CreaseGuard}{Python, PyTorch, OpenCV, FastAPI, ONNX}{In progress}
\begin{itemize}
  \item Built a cricket no-ball detection system: two residual CNNs with squeeze-and-excitation attention classify front-foot oversteps and waist-high full tosses from side-on camera frames, with Grad-CAM heatmaps to make each decision checkable
  \item Added a rule engine for the remaining no-ball laws across five match formats and a FastAPI service that evaluates new frames automatically; live broadcast detection is the next stage
\end{itemize}
\vspace{2pt}
\project{AML Pattern Engine}{C++20, pybind11, Python}{Design stage}
\begin{itemize}
  \item Designing a temporal graph engine to detect money-laundering patterns (cycles, fan-in, fan-out, scatter-gather) in IBM's synthetic dataset of 5.08 million transactions; design document, C++ API and test fixtures complete
\end{itemize}
\vspace{2pt}
\project{Trustworthy Sensor Node}{Bare-metal C, FastAPI}{Design stage}
\begin{itemize}
  \item Designed a packet format for HMAC-SHA256-signed, replay-protected sensor readings from a microcontroller to a gateway that detects forgery, replay and tampering
\end{itemize}
\vspace{2pt}
\project{GradPilot}{Team hackathon project}{State-level finalist, 2025}
\begin{itemize}
  \item Co-developed an AI admissions planner that builds a weekly-adapting roadmap from past admits to a student's target universities; state-level finalist at Bit N Build 2025 (Google Developer Groups)
\end{itemize}

\section{Leadership}
\entry{Business Innovators Community, VIT Chennai}{July 2026 -- Present}
\role{Research and Development Lead}{Chennai, India}
\begin{itemize}
  \item Lead an 8-member research and development team; edit and review articles for its Medium account (2 published)
  \item Supported sponsorship outreach and the organisation of DeFy, the community's two-day Web3 hackathon
\end{itemize}

\end{document}
"
\providecommand{\phoneline}{}

\begin{document}

\begin{center}
  {\LARGE\bfseries Cedric Antony Joseph}\\[5pt]
  Chennai, India \textbar{} \phoneline
  \href{mailto:cedricaj31@gmail.com}{cedricaj31@gmail.com}\\[2pt]
  \href{https://www.linkedin.com/in/cedric-antony-joseph-5b68b4224/}{linkedin.com/in/cedric-antony-joseph-5b68b4224} \textbar{}
  \href{https://github.com/Divide317}{github.com/Divide317}
\end{center}
\vspace{-6pt}

\section{Summary}
Computer science undergraduate moving from AI and machine learning into cybersecurity and systems. Building detection engines at three layers of the stack: kernel (Linux rootkit detector), device (signed sensor telemetry) and ledger (money-laundering engine).

\section{Education}
\entry{Vellore Institute of Technology (VIT), Chennai}{Expected June 2028}
\role{B.Tech in Computer Science and Engineering (Specialisation in Artificial Intelligence and Machine Learning)}{Chennai, India}
\begin{itemize}
  \item \textbf{Relevant Coursework:} Operating Systems, Computer Networks, Data Structures and Algorithms, Database Management Systems, Machine Learning, Probability and Statistics
  \item \textbf{Certification in progress:} Google Cybersecurity Professional Certificate (Coursera)
\end{itemize}

\section{Technical Skills}
\textbf{Languages:} C, C++, Python, Java, SQL, JavaScript\par
\vspace{1.5pt}
\textbf{Systems and Security:} Linux internals (/proc, syscalls), rootkit detection, MITRE ATT\&CK, threat modelling, Kali Linux\par
\vspace{1.5pt}
\textbf{Security Operations (certificate in progress):} SIEM tools, intrusion detection systems, packet analysis, incident response\par
\vspace{1.5pt}
\textbf{Testing and Tooling:} Git, GitHub Actions, AddressSanitizer, UBSan, Valgrind, clang-tidy, pytest, Docker\par
\vspace{1.5pt}
\textbf{Machine Learning and Backend:} PyTorch, TensorFlow, scikit-learn, OpenCV, ONNX, FastAPI, PostgreSQL, MongoDB, AWS\par

\section{Experience}
\entry{Superposition Technologies}{One-month internship}
\role{AI Intern}{}
\begin{itemize}
  \item Contributed to an AI-driven platform that routes inbound calls to the right recipient and schedules recurring weekly meetings without manual coordination
\end{itemize}

\section{Projects}
\project{Crossview}{C, Python, GitHub Actions}{In progress}
\begin{itemize}
  \item Building a Linux rootkit detector: a C library asks the kernel about processes, modules, ports and preload hooks through independent interfaces (libc, raw syscalls, direct probes) and flags where the answers disagree
  \item Localises each mismatch to user-space or kernel-level hiding and maps it to a MITRE ATT\&CK technique; Python command-line tool over a ctypes interface with JSON reports
  \item CI runs the tests under AddressSanitizer, UBSan and Valgrind with clang-tidy and pytest; published threat model and limitations
\end{itemize}
\vspace{2pt}
\project{CreaseGuard}{Python, PyTorch, OpenCV, FastAPI, ONNX}{In progress}
\begin{itemize}
  \item Built a cricket no-ball detection system: two residual CNNs with squeeze-and-excitation attention classify front-foot oversteps and waist-high full tosses from side-on camera frames, with Grad-CAM heatmaps to make each decision checkable
  \item Added a rule engine for the remaining no-ball laws across five match formats and a FastAPI service that evaluates new frames automatically; live broadcast detection is the next stage
\end{itemize}
\vspace{2pt}
\project{AML Pattern Engine}{C++20, pybind11, Python}{Design stage}
\begin{itemize}
  \item Designing a temporal graph engine to detect money-laundering patterns (cycles, fan-in, fan-out, scatter-gather) in IBM's synthetic dataset of 5.08 million transactions; design document, C++ API and test fixtures complete
\end{itemize}
\vspace{2pt}
\project{Trustworthy Sensor Node}{Bare-metal C, FastAPI}{Design stage}
\begin{itemize}
  \item Designed a packet format for HMAC-SHA256-signed, replay-protected sensor readings from a microcontroller to a gateway that detects forgery, replay and tampering
\end{itemize}
\vspace{2pt}
\project{GradPilot}{Team hackathon project}{State-level finalist, 2025}
\begin{itemize}
  \item Co-developed an AI admissions planner that builds a weekly-adapting roadmap from past admits to a student's target universities; state-level finalist at Bit N Build 2025 (Google Developer Groups)
\end{itemize}

\section{Leadership}
\entry{Business Innovators Community, VIT Chennai}{July 2026 -- Present}
\role{Research and Development Lead}{Chennai, India}
\begin{itemize}
  \item Lead an 8-member research and development team; edit and review articles for its Medium account (2 published)
  \item Supported sponsorship outreach and the organisation of DeFy, the community's two-day Web3 hackathon
\end{itemize}

\end{document}
"
\providecommand{\phoneline}{}

\begin{document}

\begin{center}
  {\LARGE\bfseries Cedric Antony Joseph}\\[5pt]
  Chennai, India \textbar{} \phoneline
  \href{mailto:cedricaj31@gmail.com}{cedricaj31@gmail.com}\\[2pt]
  \href{https://www.linkedin.com/in/cedric-antony-joseph-5b68b4224/}{linkedin.com/in/cedric-antony-joseph-5b68b4224} \textbar{}
  \href{https://github.com/Divide317}{github.com/Divide317}
\end{center}
\vspace{-6pt}

\section{Summary}
Computer science undergraduate moving from AI and machine learning into cybersecurity and systems. Building detection engines at three layers of the stack: kernel (Linux rootkit detector), device (signed sensor telemetry) and ledger (money-laundering engine).

\section{Education}
\entry{Vellore Institute of Technology (VIT), Chennai}{Expected June 2028}
\role{B.Tech in Computer Science and Engineering (Specialisation in Artificial Intelligence and Machine Learning)}{Chennai, India}
\begin{itemize}
  \item \textbf{Relevant Coursework:} Operating Systems, Computer Networks, Data Structures and Algorithms, Database Management Systems, Machine Learning, Probability and Statistics
  \item \textbf{Certification in progress:} Google Cybersecurity Professional Certificate (Coursera)
\end{itemize}

\section{Technical Skills}
\textbf{Languages:} C, C++, Python, Java, SQL, JavaScript\par
\vspace{1.5pt}
\textbf{Systems and Security:} Linux internals (/proc, syscalls), rootkit detection, MITRE ATT\&CK, threat modelling, Kali Linux\par
\vspace{1.5pt}
\textbf{Security Operations (certificate in progress):} SIEM tools, intrusion detection systems, packet analysis, incident response\par
\vspace{1.5pt}
\textbf{Testing and Tooling:} Git, GitHub Actions, AddressSanitizer, UBSan, Valgrind, clang-tidy, pytest, Docker\par
\vspace{1.5pt}
\textbf{Machine Learning and Backend:} PyTorch, TensorFlow, scikit-learn, OpenCV, ONNX, FastAPI, PostgreSQL, MongoDB, AWS\par

\section{Experience}
\entry{Superposition Technologies}{One-month internship}
\role{AI Intern}{}
\begin{itemize}
  \item Contributed to an AI-driven platform that routes inbound calls to the right recipient and schedules recurring weekly meetings without manual coordination
\end{itemize}

\section{Projects}
\project{Crossview}{C, Python, GitHub Actions}{In progress}
\begin{itemize}
  \item Building a Linux rootkit detector: a C library asks the kernel about processes, modules, ports and preload hooks through independent interfaces (libc, raw syscalls, direct probes) and flags where the answers disagree
  \item Localises each mismatch to user-space or kernel-level hiding and maps it to a MITRE ATT\&CK technique; Python command-line tool over a ctypes interface with JSON reports
  \item CI runs the tests under AddressSanitizer, UBSan and Valgrind with clang-tidy and pytest; published threat model and limitations
\end{itemize}
\vspace{2pt}
\project{CreaseGuard}{Python, PyTorch, OpenCV, FastAPI, ONNX}{In progress}
\begin{itemize}
  \item Built a cricket no-ball detection system: two residual CNNs with squeeze-and-excitation attention classify front-foot oversteps and waist-high full tosses from side-on camera frames, with Grad-CAM heatmaps to make each decision checkable
  \item Added a rule engine for the remaining no-ball laws across five match formats and a FastAPI service that evaluates new frames automatically; live broadcast detection is the next stage
\end{itemize}
\vspace{2pt}
\project{AML Pattern Engine}{C++20, pybind11, Python}{Design stage}
\begin{itemize}
  \item Designing a temporal graph engine to detect money-laundering patterns (cycles, fan-in, fan-out, scatter-gather) in IBM's synthetic dataset of 5.08 million transactions; design document, C++ API and test fixtures complete
\end{itemize}
\vspace{2pt}
\project{Trustworthy Sensor Node}{Bare-metal C, FastAPI}{Design stage}
\begin{itemize}
  \item Designed a packet format for HMAC-SHA256-signed, replay-protected sensor readings from a microcontroller to a gateway that detects forgery, replay and tampering
\end{itemize}
\vspace{2pt}
\project{GradPilot}{Team hackathon project}{State-level finalist, 2025}
\begin{itemize}
  \item Co-developed an AI admissions planner that builds a weekly-adapting roadmap from past admits to a student's target universities; state-level finalist at Bit N Build 2025 (Google Developer Groups)
\end{itemize}

\section{Leadership}
\entry{Business Innovators Community, VIT Chennai}{July 2026 -- Present}
\role{Research and Development Lead}{Chennai, India}
\begin{itemize}
  \item Lead an 8-member research and development team; edit and review articles for its Medium account (2 published)
  \item Supported sponsorship outreach and the organisation of DeFy, the community's two-day Web3 hackathon
\end{itemize}

\end{document}
"
\providecommand{\phoneline}{}

\begin{document}

\begin{center}
  {\LARGE\bfseries Cedric Antony Joseph}\\[5pt]
  Chennai, India \textbar{} \phoneline
  \href{mailto:cedricaj31@gmail.com}{cedricaj31@gmail.com}\\[2pt]
  \href{https://www.linkedin.com/in/cedric-antony-joseph-5b68b4224/}{linkedin.com/in/cedric-antony-joseph-5b68b4224} \textbar{}
  \href{https://github.com/Divide317}{github.com/Divide317}
\end{center}
\vspace{-6pt}

\section{Summary}
Computer science undergraduate moving from AI and machine learning into cybersecurity and systems. Building detection engines at three layers of the stack: kernel (Linux rootkit detector), device (signed sensor telemetry) and ledger (money-laundering engine).

\section{Education}
\entry{Vellore Institute of Technology (VIT), Chennai}{Expected June 2028}
\role{B.Tech in Computer Science and Engineering (Specialisation in Artificial Intelligence and Machine Learning)}{Chennai, India}
\begin{itemize}
  \item \textbf{Relevant Coursework:} Operating Systems, Computer Networks, Data Structures and Algorithms, Database Management Systems, Machine Learning, Probability and Statistics
  \item \textbf{Certification in progress:} Google Cybersecurity Professional Certificate (Coursera)
\end{itemize}

\section{Technical Skills}
\textbf{Languages:} C, C++, Python, Java, SQL, JavaScript\par
\vspace{1.5pt}
\textbf{Systems and Security:} Linux internals (/proc, syscalls), rootkit detection, MITRE ATT\&CK, threat modelling, Kali Linux\par
\vspace{1.5pt}
\textbf{Security Operations (certificate in progress):} SIEM tools, intrusion detection systems, packet analysis, incident response\par
\vspace{1.5pt}
\textbf{Testing and Tooling:} Git, GitHub Actions, AddressSanitizer, UBSan, Valgrind, clang-tidy, pytest, Docker\par
\vspace{1.5pt}
\textbf{Machine Learning and Backend:} PyTorch, TensorFlow, scikit-learn, OpenCV, ONNX, FastAPI, PostgreSQL, MongoDB, AWS\par

\section{Experience}
\entry{Superposition Technologies}{One-month internship}
\role{AI Intern}{}
\begin{itemize}
  \item Contributed to an AI-driven platform that routes inbound calls to the right recipient and schedules recurring weekly meetings without manual coordination
\end{itemize}

\section{Projects}
\project{Crossview}{C, Python, GitHub Actions}{In progress}
\begin{itemize}
  \item Building a Linux rootkit detector: a C library asks the kernel about processes, modules, ports and preload hooks through independent interfaces (libc, raw syscalls, direct probes) and flags where the answers disagree
  \item Localises each mismatch to user-space or kernel-level hiding and maps it to a MITRE ATT\&CK technique; Python command-line tool over a ctypes interface with JSON reports
  \item CI runs the tests under AddressSanitizer, UBSan and Valgrind with clang-tidy and pytest; published threat model and limitations
\end{itemize}
\vspace{2pt}
\project{CreaseGuard}{Python, PyTorch, OpenCV, FastAPI, ONNX}{In progress}
\begin{itemize}
  \item Built a cricket no-ball detection system: two residual CNNs with squeeze-and-excitation attention classify front-foot oversteps and waist-high full tosses from side-on camera frames, with Grad-CAM heatmaps to make each decision checkable
  \item Added a rule engine for the remaining no-ball laws across five match formats and a FastAPI service that evaluates new frames automatically; live broadcast detection is the next stage
\end{itemize}
\vspace{2pt}
\project{AML Pattern Engine}{C++20, pybind11, Python}{Design stage}
\begin{itemize}
  \item Designing a temporal graph engine to detect money-laundering patterns (cycles, fan-in, fan-out, scatter-gather) in IBM's synthetic dataset of 5.08 million transactions; design document, C++ API and test fixtures complete
\end{itemize}
\vspace{2pt}
\project{Trustworthy Sensor Node}{Bare-metal C, FastAPI}{Design stage}
\begin{itemize}
  \item Designed a packet format for HMAC-SHA256-signed, replay-protected sensor readings from a microcontroller to a gateway that detects forgery, replay and tampering
\end{itemize}
\vspace{2pt}
\project{GradPilot}{Team hackathon project}{State-level finalist, 2025}
\begin{itemize}
  \item Co-developed an AI admissions planner that builds a weekly-adapting roadmap from past admits to a student's target universities; state-level finalist at Bit N Build 2025 (Google Developer Groups)
\end{itemize}

\section{Leadership}
\entry{Business Innovators Community, VIT Chennai}{July 2026 -- Present}
\role{Research and Development Lead}{Chennai, India}
\begin{itemize}
  \item Lead an 8-member research and development team; edit and review articles for its Medium account (2 published)
  \item Supported sponsorship outreach and the organisation of DeFy, the community's two-day Web3 hackathon
\end{itemize}

\end{document}
